\documentclass[11pt,a4paper]{article}

\usepackage[utf8]{inputenc}
\usepackage[T1]{fontenc}
\usepackage[margin=2.5cm]{geometry}
\usepackage{amsmath,amssymb}
\usepackage{booktabs}
\usepackage{tabularx}
\usepackage{algorithm}
\usepackage{algpseudocode}
\usepackage{enumitem}
\usepackage[hidelinks]{hyperref}

\newcommand{\code}[1]{\texttt{#1}}
\newcommand{\cand}{\mathcal{C}}
\newcommand{\key}{\kappa}
\algrenewcommand\algorithmicrequire{\textbf{Input:}}
\algrenewcommand\algorithmicensure{\textbf{Output:}}

\usepackage[backend=biber, style=apa]{biblatex}
\addbibresource{references.bib}


\title{Choice-Set Generation for Public Transport Route Choice\\
\large Algorithm of \code{run\_choice\_set.py}}
\author{Símun Pauli Pálsson}
\date{\today}

\begin{document}
\maketitle

\section{Overview}

For every trip in the Danish National Travel Survey (TU) the algorithm builds a \emph{choice set}: the
chosen itinerary plus a set of distinct public transport alternatives between the same origin and destination around the same departure time. Alternatives are generated with OpenTripPlanner (OTP) by repeatedly querying it under perturbed conditions --- stops and routes elimination, random generalized-cost parameters, and the fastest routes successively removed --- and then collapsing the pooled itineraries to one representative per distinct alternative.

Two choice sets are produced:
\begin{description}[leftmargin=1.5em]
  \item[Public transport trips.] The chosen itinerary is the OTP itinerary matched to the TU trip by
    the reconstruction run (\code{run\_reconstruct.py}). Only trips whose legs are all walk or transit
    are used, and trips found only by ignoring the reported route (\code{trip\_wrong\_route}) are
    dropped, since their route is not the chosen one.
  \item[Car trips.] Trips with car as primary mode and only walk/car legs. The chosen alternative is
    the TU car trip itself (distance and duration as reported); the alternatives are the public
    transport itineraries the respondent could have taken instead, with walk access/egress.
\end{description}

The alternative generation code (\code{src/choice\_set/}) takes plain data --- a query, OTP itineraries, a
departure time, an optional chosen itinerary --- and does \textbf{not depend on TU}, so it can be reused for
other travel surveys.

\paragraph{Notation.}
For a trip $i$ let $o_i, d_i$ be origin and destination, $t_i$ the reported departure time, $W$ the
search window (\code{choice\_set.search\_window}, e.g.\ 1 hour) and $c_i$ the chosen itinerary. An
itinerary $r$ is a sequence of legs; $T(r)$ denotes its transit legs in order, $s(r)$ and $e(r)$ its
start and end time. $\cand$ denotes a pool of candidate itineraries.



\section{OTP Routing Algorithm}
OTP uses the McRAPTOR algorithm \parencite{RAPTOR}, a round-based public transit router, which give a pareto-optimal set of itineries for multi-critera. OTP allows for these criterias: arrival, in-vechical time, and number of transfers. This gives a much faster algorithm then traditional Dijkstra shortest-path algorithms where transit schedule is build into the graph.

To find all stops within reach (a predetermined duration) of origin and destination, the algortihm starts by expanding street mode in all direction (using Dijkstra) from origin and destination.
%This section needs more work
%Explain Query $Q$



\section{The base OTP query}
\label{sec:query}

Each OTP request, $Q(\theta)$, plans transit itineraries from $o_i$ to $d_i$ departing at $t_i$ with
the trip's access and egress mode, on \emph{all} transit modes (bus, rail, S-train, metro, tram) and
all routes. $\theta$ collects the optional perturbations of Section~\ref{sec:rounds}: excluded stops,
excluded routes, and cost parameters (walk, wait transfer, and per-mode transit cost);
without them OTP's defaults apply (walk cost 2).

\paragraph{Pagination.} A single OTP response holds only a page of itineraries, so $Q$ pages forward
from $t_i$ until the latest itinerary starts after $t_i + W$, and backward until the earliest starts
before $t_i - W$, at most \code{max\_itinerary\_candidates} itineraries each way. Duplicated itineraries are removed (TODO: define duplication).

\paragraph{Direct-walk filter.} Let $D_i$ be the distance of the direct walk from $o_i$ to $d_i$. An
itinerary is dropped if its walking legs together are at least as long as the direct walk:
\begin{equation}
  \text{keep } r \iff \sum_{\ell \in r,\ \text{mode}(\ell)=\text{WALK}} \text{dist}(\ell) + 0.01\,\text{km} < D_i .
\end{equation}
Such itineraries use transit without saving any walking and are not meaningful alternatives. This is
clearly a conservative threshold. OTP has many more itinerary filters, but they are designed to show
users only ``good'' itineraries. In the TU data some respondents go to great lengths to use transit,
leaving one to wonder how that choice is better than simply walking directly. This can, however, be
explained by misjudgement, habit, or bad weather, which may lead respondents to disregard the cost of
walking directly and consider only transit route.

%\paragraph{Bicycle access/egress.} OTP's bicycle mode requires both access and egress by bicycle, so
%TU bicycle access/egress is queried as walking. After generation, the first (last) walk leg of each
%itinerary is relabelled \code{BICYCLE} and its duration multiplied by
%\code{tu\_prep.walk\_bike\_time\_ratio}, as in the reconstruction.

\section{Generating candidates}
\label{sec:rounds}

Candidate generation starts from the plain query $\cand_0 = Q(\emptyset)$ and then runs a fixed,
ordered sequence of \emph{rounds}. Each round issues one re-query $Q(\theta_k)$ and appends its
itineraries to the pool, tagged with the round number and its perturbation. Itineraries found in
several rounds are kept in the pool at this stage; duplicates are collapsed in
Section~\ref{sec:collapse}. The round types run in this order:

\begin{enumerate}
  \item stop exclusion rounds (\code{choice\_set.exclude\_stops}),
  \item route exclusion rounds (\code{choice\_set.exclude\_routes}),
  \item random cost rounds (\code{choice\_set.random\_costs}),
  \item targeted dominating route exclusion rounds (\code{choice\_set.exclude\_targeted}).
\end{enumerate}

Targeted rounds run last, which means the routes and stops they discover do not enter the random route/stop exclusion rounds.



\subsection{Random stop and route exclusion}

Exclusion forces OTP away from the stops and routes it otherwise prefers, revealing alternatives that
are close in cost but never optimal. The exclusion rounds are organized in \emph{steps}, each a number
of rounds excluding $n$ ids per round (e.g.\ 5 rounds of 1 stop, then 5 rounds of 2 stops).

In round $k$ with $n$ ids to exclude, the id pool $P_k$ is the set of ids used by the transit legs of
the chosen itinerary and of all candidates found so far:
\begin{itemize}
  \item \emph{stops}: boarding and alighting stop ids. An excluded stop cannot be used for boarding or
    alighting, but trips may still pass through it;
  \item \emph{routes}: route ids. An excluded route's trips are not used at all.
\end{itemize}
Because the pool grows with every round, later rounds can exclude stops or routes of alternatives
that earlier rounds uncovered. The excluded set is drawn uniformly from the $\min(n, |P_k|)$-subsets
of $P_k$ that no earlier round of the same kind has excluded:
\begin{equation}
  E_k \sim \text{Uniform} \Bigl(\bigl\{ S \subseteq P_k \setminus \mathcal{E}: |S| = \min(n,|P_k|\setminus \mathcal{E}) \bigr\} \Bigr),
\end{equation}
where $\mathcal{E}$ is the set of all sets already excluded (across steps). If no such set is left,
the round is skipped.


\subsection{Random generalized costs}

Each cost round queries without exclusions but with cost parameters drawn independently per round:
\begin{align}
  \beta_{\text{walk}} &\sim \text{LogUniform}[a_{\text{walk}}, b_{\text{walk}}], &
  \beta_{\text{wait}} &\sim \text{LogUniform}[a_{\text{wait}}, b_{\text{wait}}], \\
  \beta_{m} &\sim \text{LogUniform}[a_{\text{transit}}, b_{\text{transit}}] \quad \forall m \in \mathcal{M}, &
  c_{\text{transfer}} &\sim \text{Uniform}[a_{\text{tr}}, b_{\text{tr}}],
\end{align}
where $\beta_{\text{walk}}$ and $\beta_{\text{wait}}$ are OTP's walk and wait cost, $\beta_m$ the in-vehicle cost of transit mode $m \in \mathcal{M} = \{\text{bus}, \text{rail}, \text{S-train}, \text{metro}, \text{tram}\}$, and $c_{\text{transfer}}$ a cost (in seconds) added to every boarding after the first. Costs are rounded to $0.01$, the transfer cost to whole seconds.
A draw $X \sim \text{LogUniform}[a,b]$ is $X = \exp(U)$ with $U \sim \mathcal{U}[\ln a, \ln b]$, so that, e.g., halving and doubling a cost are equally likely. Varying these weights mimics travellers with different preferences (e.g.\ walk-averse, transfer-averse, or with a preference for rail), each yielding the itineraries optimal for that preference.

\subsection{Targeted exclusion of the fastest routes}
Random exclusion may not remove the right combination of routes when one fast route dominates
another, e.g.\ a bus running along a metro line. Targeted rounds peel off the fastest routes
deterministically. Let $\mathcal{F}_j$ be the itineraries found by targeted round $j$
($\mathcal{F}_0 = \cand_0$), $E^*_j$ the routes it excludes ($E^*_0 = \emptyset$) and $R(r)$ the routes
of itinerary $r$. Round $j$ adds the routes of the previous round's fastest itinerary $r^*_j$ (ties to
the earlier start) to those already excluded:
\begin{equation}
  r^*_j = \operatorname*{arg\,min}_{r \in \mathcal{F}_{j-1}} \bigl(e(r) - s(r)\bigr), \qquad
  E^*_j = E^*_{j-1} \cup R(r^*_j), \qquad
  \mathcal{F}_j = Q(\text{exclude routes } E^*_j).
\end{equation}
Once a round finds no itinerary, the remaining targeted rounds are skipped.
For example, if the fastest itinerary of the plain query is metro M1, round 1 excludes M1 and may find bus 5C along the same corridor; round 2 then excludes both M1 and 5C, and so on.

\subsection{Reproducibility}

Every random draw uses a random number generator seeded per trip from the configured \code{seed} and the trip's \code{TurId}, with a separate generator for stops, routes and costs. The draws thus do not depend on the order in which the parallel worker threads process trips, and repeated runs produce the same choice sets (given the same OTP graph).

\begin{algorithm}[htbp]
\caption{Candidate generation for one trip $i$}
\label{alg:candidates}
\begin{algorithmic}[1]
\Require query $Q$ (with pagination and direct-walk filter), chosen itinerary $c_i$, rounds $R = (\rho_1,\dots,\rho_K)$
\Ensure candidate pool $\cand$
\State $\cand \gets Q(\emptyset)$; \quad $\cand^{\text{t}} \gets \cand$; \quad $X \gets \emptyset$; \quad $\mathcal{E} \gets \emptyset$
\For{$k = 1, \dots, K$}
  \If{$\rho_k$ is a cost round}
    \State draw $\theta_k = (\beta_{\text{walk}}, \beta_{\text{wait}}, (\beta_m)_{m}, c_{\text{transfer}})$
  \ElsIf{$\rho_k$ is a targeted round}
    \State $r^* \gets$ fastest itinerary of $\cand^{\text{t}}$
    \If{$r^*$ does not exist} \textbf{continue} \EndIf
    \State $X \gets X \cup \text{routes}(T(r^*))$; \quad $\theta_k \gets (\text{exclude routes } X)$
  \Else \Comment{stop or route exclusion with $n$ ids}
    \State $P \gets$ stop (route) ids of the transit legs of $c_i$ and $\cand$
    \State $\mathcal{S} \gets \{ S \subseteq P : |S| = \min(n,|P|),\ S \notin \mathcal{E} \}$
    \If{$\mathcal{S} = \emptyset$} \textbf{continue} \EndIf
    \State draw $E$ uniformly from $\mathcal{S}$; \quad $\mathcal{E} \gets \mathcal{E} \cup \{E\}$; \quad $\theta_k \gets (\text{exclude } E)$
  \EndIf
  \State $\cand_k \gets Q(\theta_k)$
  \If{$\rho_k$ is a targeted round} $\cand^{\text{t}} \gets \cand_k$ \EndIf
  \State $\cand \gets \cand \cup \cand_k$ \Comment{tagged with round $k$ and $\theta_k$}
\EndFor
\State relabel bicycle access/egress (Section~\ref{sec:query})
\State \Return $\cand$
\end{algorithmic}
\end{algorithm}

\section{Collapsing candidates into alternatives}
\label{sec:collapse}

The pool contains many itineraries that are, from the route choice point of view, the same
alternative --- the same routes departing at different times, or the same itinerary found in several
rounds. The pool is collapsed to one itinerary per distinct alternative.

\paragraph{Alternative key.} Two itineraries are the same alternative if they have the same key
$\key(r)$, a tuple over their transit legs in order. The key is configured by
\code{choice\_set.alternative\_key}:

\begin{center}
\begin{tabularx}{\linewidth}{@{}lX@{}}
\toprule
\code{alternative\_key} & $\key(r)$, per transit leg $\ell \in T(r)$ \\
\midrule
\code{routes} & (mode, route name) \\
\code{routes\_stops} & (mode, route name, boarding stop name, alighting stop name) \\
\code{trip} & (mode, route name, boarding stop name, alighting stop name, departure time) \\
\bottomrule
\end{tabularx}
\end{center}

Stops are compared by name rather than by stop id, so the same route boarded at another platform of
the same station is not a new alternative. With \code{trip}, every departure is its own alternative.

\paragraph{Representative itinerary.} Among the candidates sharing a key, the itinerary kept is the
one closest to the chosen departure, with waiting counted as a loss. With $t^{c}_i$ the start of the
chosen itinerary (public transport trips) or the reported departure time (car trips), each
candidate's score in minutes is
\begin{equation}
  \text{score}(r) = \bigl| s(r) - t^{c}_i \bigr| + w(r),
  \qquad
  w(r) = \sum_{j=2}^{|r|} \bigl( s(\ell_j) - e(\ell_{j-1}) \bigr),
\end{equation}
where $w(r)$ is the total waiting time between consecutive legs. The candidate with the lowest score is
kept, ties going to the one with less waiting. This prefers, of two itineraries arriving at the same
time, the one leaving later rather than the one leaving earlier and waiting at a transfer.

\paragraph{Chosen itinerary.} The chosen itinerary $c_i$ is alternative $0$. For public transport
trips, candidates with the chosen key $\key(c_i)$ are removed: the chosen itinerary itself represents
its alternative. The remaining representatives are numbered $1, 2, \dots$ in increasing score.

\begin{algorithm}[htbp]
\caption{Building the choice set of trip $i$}
\label{alg:collapse}
\begin{algorithmic}[1]
\Require candidate pool $\cand$, chosen itinerary $c_i$, reference time $t^{c}_i$, key function $\key$
\Ensure choice set $A_i = (a_0, a_1, \dots)$
\State $a_0 \gets c_i$
\If{public transport trip}
  \State $\text{found} \gets \exists\, r \in \cand : \key(r) = \key(c_i)$
  \State $\cand \gets \{ r \in \cand : \key(r) \neq \key(c_i) \}$
\EndIf
\State sort $\cand$ by $(\text{score}(r), w(r))$ ascending
\State keep the first itinerary of each key $\to (a_1, a_2, \dots)$
\State assign line groups (below); compute coverage (public transport trips)
\State \Return $A_i$
\end{algorithmic}
\end{algorithm}

\paragraph{Line groups.} Some distinct alternatives differ only by a parallel line serving the same
stops, e.g.\ buses 143 and 144, metro M1 and M2, or S-train lines A and C on a shared section. They are
kept as separate alternatives but grouped in a \code{line\_group}: each alternative's group is the
lowest alternative number with the same \emph{line-blind} key, i.e.\ the key over (alternative-key
columns, boarding stop, alighting stop) with the route name replaced by
\begin{itemize}
  \item nothing for buses (all buses between the same stops form one group);
  \item the lower of two metro lines sharing track (M2 $\to$ M1, M4 $\to$ M3);
  \item one common name for all S-train lines except F (which does not share track with the others);
  \item the route name itself for other modes.
\end{itemize}
Line groups let a model treat overlapping alternatives (e.g.\ with a path-size correction) or
aggregate them.

\paragraph{Coverage.} For public transport trips, the choice set is \emph{covered} if the generation
reproduced the chosen alternative, i.e.\ if some candidate had the chosen key, or some generated
alternative falls in the chosen itinerary's line group (line group $0$):
\begin{equation}
  \text{coverage}_i = \mathbb{1}\Bigl[\, \exists\, r \in \cand : \key(r) = \key(c_i)
  \ \lor\ \exists\, a_j \in A_i,\ j>0 : \text{line\_group}(a_j) = 0 \Bigr].
\end{equation}
Coverage measures how well the generation procedure spans the travellers' actual choices; uncovered
trips point to alternatives the procedure cannot generate.

\section{Output}

Each choice set is written with one row per leg, identified by \code{TurId} and \code{alt\_id}
($0$ = chosen, with \code{is\_chosen} $=1$), carrying the leg's mode, route, stops, times, distance,
straight-line distance, OTP generalized cost and geometry, plus \code{line\_group} and (for public
transport trips) \code{coverage}. A summary reports per choice set (public / car) and, for public
transport trips, per reconstruction outcome: the number of trips, the number and share covered, the
number of trips without any OTP transit itinerary, and the mean and percentiles of the number of
alternatives besides the chosen.

The procedure deliberately over-generates: the choice set may contain many alternatives, including
ones a traveller would rarely consider.

\section{Configuration}

\begin{center}
\begin{tabularx}{\linewidth}{@{}lX@{}}
\toprule
Key (\code{choice\_set.*}) & Meaning \\
\midrule
\code{reconstruct\_run\_dir} & Reconstruction run whose matched itineraries are the chosen trips
  (\code{null}: newest run of the year) \\
\code{search\_window} & $W$: alternatives start within $t_i \pm W$ (ISO 8601 duration) \\
\code{alternative\_key} & \code{routes}, \code{routes\_stops} or \code{trip} (Section~\ref{sec:collapse}) \\
\code{exclude\_stops} & \code{steps} of (\code{rounds}, \code{stops\_per\_round}); \code{seed} \\
\code{exclude\_routes} & \code{steps} of (\code{rounds}, \code{routes\_per\_round}); \code{seed} \\
\code{random\_costs} & \code{rounds}; $[\min,\max]$ of \code{walk\_cost}, \code{wait\_cost},
  \code{transit\_cost}, \code{transfer\_cost}; \code{seed} \\
\code{exclude\_targeted} & \code{rounds} \\
\bottomrule
\end{tabularx}
\end{center}

With the example configuration (\code{README.md}), each trip is queried $1 + 10 + 10 + 10 + 4 = 35$
times: the plain query, 5 rounds excluding one and 5 excluding two stops, likewise for routes, 10
random cost rounds and 4 targeted rounds, each searching $\pm 1$ hour around the departure.

\end{document}